\subsection{Convex cones with compact sole}

PROPOSITION 6. --- Let E be a Hausdorff locally convex space and K a convex compact set in E which does not contain 0. Then the smallest pointed cone C of vertex 0 which contains K is a proper convex cone in E and is a locally compact and complete subspace of E; also, there exists a closed hyperplane H in E that does not contain 0 and is such that H meets all the half-lines originating at 0 contained in C and such that H ∩ C is compact. Further, if D is the half-space containing 0 determined by H, a closed hyperplane with these properties, then C ∩ D is a cap of C and C is the union of the λ(C ∩ D) for λ \textgreater{} 0.

By prop. 4 of II, p.~38, there exists a closed hyperplane H which separates 0 strictly from K. Now, the convex envelope A of the union of \(\{0\}\) and of K is compact (II, p.~14, prop. 15) and is the union of the \(\lambda K\) with \(0 \leqslant \lambda \leqslant 1\) . As 0 and K are strictly on opposite sides of H, for every \(x \in K\) there exists \(\lambda\) such that \(0 < \lambda < 1\) and \(\lambda x \in H\) . As C is the union of the \(\lambda A\) for \(\lambda \geqslant 1\) , we see that H meets every half-line originating at 0 contained in C and that \(H \cap A = H \cap C\) is compact. Further, C is also the union of the \(\lambda(H \cap C)\) for \(\lambda \geqslant 0\) ; let \(C_n\) be the union of the \(\lambda(H \cap C)\) for \(0 \leqslant \lambda \leqslant n\) . Clearly \(C_n\) is the convex envelope of the union of \(\{0\}\) and of \(n(H \cap C)\) , therefore it is compact. Also, for all \(x \in E\) , there is a closed neighbourhood V of x in E and an integer n such that \(V \cap C \subset C_n\) ; in fact, if H is defined by the equation \(f(z) = \alpha\) , where \(\alpha > 0\) , it is sufficient to take for V the closed half-space determined by nH and containing 0, where n is so large that \(n\alpha > f(x)\) . This shows that C is locally compact (taking \(x \in C\) ), and that it is closed in E. We can also consider K as a subset of the completion \(\hat{E}\) , therefore C is also closed in \(\hat{E}\) and therefore complete.

Given a cone C and a closed hyperplane H in a Hausdorff topological vector space E, such that H does not contain the vertex s of C and C is the smallest cone with vertex s containing \(H \cap C\) , then we call the intersection \(H \cap C\) a « sole » of the cone C. Prop. 6 shows that in a Hausdorff locally convex space E, the smallest cone of vertex 0, containing a compact convex set K to which 0 does not belong, is a cone of compact sole, and that every convex cone having a compact sole S, is locally compact and complete.

Examples. --- 1) Every proper closed convex cone in E, a vector space of finite dimension, has a compact sole. In fact, by II, p.~52, prop. 11 we need only consider the case where \(E = R^{n}\) and \(C = R_{+}^{n}\) . If \((e_{i})_{1 \leqslant i \leqslant n}\) is the canonical basis of \(R^{n}\) , it is clear that the compact convex set which is the convex envelope of the \(e_{i}\) ( \(1 \leqslant i \leqslant n\) ) is a compact sole for \(R_{+}^{n}\) .

* 2) If X is a compact space, then the cone \(\mathcal{M}_{+}(\mathrm{X})\) of positive measures on X, with the vague topology, is a cone with a compact sole (INT, III, 2nd ed., § 1, No.~9, cor. 3 of prop. 15). *

